\section{Structural reductions for the canonical index}\label{sec:index-reductions}

We reduce the canonical-index theorem to a particular terminal case.
First we bound the volume-form characters when the decomposition has
several factors or is purely abelian. A remaining single nonflat factor
has only rationally connected proper rational images. This bounds the
index for noncanonical quotients and for terminal models carrying an
effective divisor of intermediate Iitaka dimension. The cases left over
have irregularity zero and the fibration property needed for the scalar
bounds of Section~\ref{sec:scalar}.

All covers and all varieties in this
section are connected unless otherwise specified.  For a normal variety,
\(\Omega_X^{[j]}\) denotes the reflexive extension of \(\Omega^j_{X_{\rm reg}}\).
The order of \(K_X\sim_{\Q}0\) is the least positive integer \(r\) such that
\(rK_X\sim0\), with linear equivalence understood for integral Weil divisors.

\subsection{Finite covers and canonical characters}

We first record the relation between the canonical order and the action on a
volume form.  Suppose that \(q:R\to X\) is a finite Galois quasi-\'etale cover
of normal projective varieties, with group \(G\), and that \(K_R\sim0\).
Choose a generating reflexive top form \(\omega\) on \(R\).  Since a regular
function on \(R\) is constant, the action on this one-dimensional space has
the form
\[
                   g^*\omega=\chi(g)\omega,
             \qquad \chi:G\longrightarrow\C^*.
\]
Then
\begin{equation}\label{eq:index-character}
          mK_X\sim0\quad\Longleftrightarrow\quad\chi^m=1.
\end{equation}
Indeed, an invariant rational pluricanonical form on \(R\) descends at the
function-field level to \(X\).  Since \(q\) is \'etale at the generic points
of divisors, the descended form has neither zeros nor poles in codimension
one when the original form does not.  It therefore trivializes the
divisorial sheaf \(\OO_X(mK_X)\).  Conversely, a trivializing form downstairs
pulls back to a scalar multiple of \(\omega^{\otimes m}\), so that this
power is invariant.  In particular, bounding the order of \(\chi\) bounds
the exact linear-equivalence order, including the Cartier index.

For any klt \(X\) with \(K_X\sim_{\Q}0\), its cyclic index cover is obtained
by normalizing the cover of \(X_{\rm reg}\) defined by an \(r\)-th root of
a primitive trivialization of \(rK_X\).  Minimality of \(r\), or equivalently
the Kummer description of this cyclic extension, makes the cover connected
of degree \(r\).  It is quasi-\'etale and has a generating reflexive volume
form.  The finite discrepancy formula shows that it is klt.  Its canonical
divisor is Cartier and linearly trivial, so its ordinary discrepancies are
integers greater than \(-1\); it is consequently canonical.  We use the
finite discrepancy formula in this form throughout; see
\cite[Proposition~5.20]{KollarMori1998}.

The singular Beauville--Bogomolov decomposition theorem applied to this
cover gives a finite quasi-\'etale cover
\begin{equation}\label{eq:index-product-cover}
        P=A\times P_1\times\cdots\times P_s\longrightarrow X,
\end{equation}
where \(A\) is an abelian variety, possibly a point, and the positive
dimensional \(P_i\) are irreducible Calabi--Yau or irreducible holomorphic
symplectic varieties in the singular sense
\cite[Definition~1.4 and Theorem~1.5]{HoringPeternell2019}.
In dimension four the decomposition is already supplied by Druel
\cite[Theorem~1.2]{Druel2018}. The arbitrary-dimensional formulation
above combines the algebraic integrability and
holonomy methods discussed by H\"oring--Peternell in the introduction
to \cite{HoringPeternell2019}; in particular they identify the
contributions of Druel and of Greb--Guenancia--Kebekus
\cite[Theorem~B and Proposition~D]{GGK2019}.
We use the following precise properties.  Each \(P_i\) is normal,
projective and canonical, with \(K_{P_i}\sim0\) and dimension \(d_i\ge2\).
For every connected normal finite quasi-\'etale cover of \(P_i\), its algebra
of reflexive forms is generated by a volume form in the Calabi--Yau case,
and by a generically nondegenerate symplectic form in the symplectic case.
In particular, there are no reflexive one-forms or reflexive
\((d_i-1)\)-forms on these covers.  A one-dimensional factor is included in
\(A\).

\begin{proposition}\label{prop:index-product-case}
There is a uniform positive integer \(N_{\rm prod}\) with the following
property.  If a projective klt fourfold \(X\), with \(K_X\sim_{\Q}0\),
admits a cover as in \eqref{eq:index-product-cover} having at least two
positive dimensional factors, or having only an abelian factor, then
\(N_{\rm prod}K_X\sim0\).
\end{proposition}

\begin{proof}
\textbf{Preserving the product directions.}
First suppose that the product has at least two positive dimensional
factors.  Denote these factors by \(B_1,\ldots,B_t\), counting the abelian
factor, when present, as a single factor.  Thus \(t\le4\),
\(1\le\dim B_i\le3\), and at most one \(B_i\) is abelian.  Take the normal
Galois closure \(R\to X\) of the cover, with group \(G\).  This is still
quasi-\'etale: over the smooth locus of \(X\) the original cover is finite
\'etale by purity, and its Galois closure is finite \'etale there as well.
The map \(R\to P\) is also quasi-\'etale.  In particular, \(R\) is canonical
and \(K_R\sim0\).  Here and below purity is used in its usual form for a
finite generically \'etale map to a regular variety
\cite[Tag~0BMB]{StacksProject}.

The product directions pull back, and extend reflexively, to a splitting
\begin{equation}\label{eq:index-directions}
                     \mathcal T_R=\bigoplus_{i=1}^t\mathcal E_i.
\end{equation}
There are no homomorphisms between two distinct summands.  To check this,
choose a nonabelian factor among the corresponding two factors, say
\(B_i\), and fix general smooth points on all the other factors.  After
shrinking the complementary parameter space, the morphism
\(R\to\prod_{j\ne i}B_j\) has normal, pure-dimensional fibres.  This follows
from normality of the geometric generic fibre in characteristic zero,
generic flatness, and openness of geometric normality.  Each of its
connected components is therefore normal and finite surjective over
\(B_i\).  The codimension-two locus excluded from the quasi-\'etale map
\(R\to P\) meets a general slice in codimension at least two, so these
component covers are quasi-\'etale.

On the smooth big open of such a component, \(\mathcal E_i\) restricts to
its tangent sheaf and every complementary summand restricts to a trivial
bundle.  The component has no reflexive one-forms by the defining property
of \(B_i\).  It has no vector fields either: contraction with its generating
volume form identifies its reflexive tangent sheaf with its sheaf of
reflexive \((\dim B_i-1)\)-forms.  A homomorphism in either direction between
the two summands thus restricts to zero: the resulting forms or vector
fields on the big open extend reflexively across its codimension-two
complement.  These slices sweep a dense open
of \(R\), proving the assertion.

Consequently the projections in \eqref{eq:index-directions} are central
idempotents in the finite-dimensional algebra
\(\operatorname{End}_R(\mathcal T_R)\).  Its center is a finite-dimensional
commutative algebra.  Its primitive idempotents determine nonzero direct
summands of positive generic rank, whose ranks sum to four; there are at
most four of them.  The group \(G\) permutes these primitive idempotents.
The kernel \(G_0\) of that permutation action has index at most \(4!\) and
fixes every block projection, since each block projection is a sum of
primitive central idempotents.  It therefore preserves every
\(\mathcal E_i\).

\smallskip\noindent\textbf{Recovering the factor fields.}
Let \(F_i\) be the relative algebraic closure of \(\C(B_i)\) in \(\C(R)\),
and let \(R_i\to B_i\) be the normalization in this finite field extension.
At the generic point, the complementary directions in
\eqref{eq:index-directions} span
\(\operatorname{Der}_{\C(B_i)}\C(R)\): this follows first on the product
and then under the finite separable extension \(\C(P)\subset\C(R)\).
In characteristic zero the elements annihilated by all these derivations
are exactly those algebraic over \(\C(B_i)\).  Thus \(G_0\) preserves
\(F_i\), and acts by birational maps on \(R_i\).

The finite map \(R_i\to B_i\) is quasi-\'etale.  For completeness, the
regular extension obtained by adjoining the geometrically integral
product of the other factors is
linearly disjoint from the finite extension \(F_i/\C(B_i)\).  Hence the
function field of
\(R_i\times\prod_{j\ne i}B_j\) is an intermediate field between
\(\C(P)\) and \(\C(R)\).  This product is normal, so the normalization
property gives finite maps
\[
 R\longrightarrow R_i\times\prod_{j\ne i}B_j
       \longrightarrow P.
\]
Ramification over a prime divisor of \(B_i\) would persist after taking this
product.  Its ramification index would then divide the ramification index
at any prolongation to \(\C(R)\), contradicting quasi-\'etaleness of
\(R\to P\).  Thus \(R_i\) is projective and canonical with
\(K_{R_i}\sim0\).  Choose its generating volume \(\omega_i\) to be the
pullback of a generating volume on \(B_i\).

\smallskip\noindent\textbf{Bounding the characters in lower dimension.}
The varieties $R_i$ have dimension at most three. We now use their
finite birational actions to bound the characters on their volume lines.
Let \(J_i\) be the image of \(G_0\) in \(\Bir(R_i)\).  Resolve the graph
of its finite birational action to obtain a smooth projective
\(J_i\)-variety \(W_i\) with a birational morphism to \(R_i\).  One can
first take the closure of the graph in the product indexed by the group,
where the action permutes the factors, and then use functorial resolution
\cite[Theorem~1.1]{BierstoneMilman2008}.  Canonicality implies that the pullback
of \(\omega_i\) is a regular top form on \(W_i\).  For every positive
integer \(m\),
\begin{equation}\label{eq:index-block-sections}
          H^0(W_i,mK_{W_i})=\C\,\omega_i^{\otimes m}.
\end{equation}
Indeed, any section on \(W_i\) gives a rational pluricanonical form on
\(R_i\) with no pole at any prime divisor there.  It extends reflexively,
and \(\OO_{R_i}(mK_{R_i})\simeq\OO_{R_i}\) has only constant global
sections.  This also proves that the character in degree \(m\) is the
\(m\)-th power of the character \(\chi_i\) on the volume line.

Form the finite quotient \(q_i:W_i\to Q_i=W_i/J_i\).  Its branch boundary is
\[
             B_{Q_i}=\sum_D(1-1/e_D)D,
\]
where \(e_D\) is the ramification order along \(D\); unramified primes
have coefficient zero.  The quotient ramification formula gives
\[
                 K_{W_i}=q_i^*(K_{Q_i}+B_{Q_i}).
\]
The total divisor on the right is \(\Q\)-Cartier and the pair
\((Q_i,B_{Q_i})\) is klt.  These are the finite quotient and discrepancy
properties of a smooth variety with zero boundary
\cite[Proposition~5.20]{KollarMori1998}; \(\Q\)-Cartierness can also be
seen by descending a sufficiently divisible local canonical power, or
by taking a norm after trivializing near the finite fibre.  Since \(J_i\)
is finite, some power of \(\omega_i\) is invariant.  Galois descent of
rational forms, followed by the displayed ramification formula, gives
\[
                     \kappa(Q_i,K_{Q_i}+B_{Q_i})\ge0.
\]

The coefficients of these boundaries lie in the fixed DCC set
\(\{1-1/e:e\in\Z_{>0}\}\), and \(1\le\dim Q_i\le3\).  Effective log Iitaka
fibrations in these dimensions
\cite[Corollary~1.5]{ChenHanLiu2023} therefore give a uniform integer
\(m_*\), made common to the three possible dimensions, for which
\[
     H^0\bigl(Q_i,\OO_{Q_i}(\lfloor m_*(K_{Q_i}+B_{Q_i})\rfloor)\bigr)
                    \ne0.
\]
A section here pulls back as a rational \(m_*\)-canonical form on
\(W_i\).  The ramification formula and coefficientwise effectiveness
show that it has no divisorial poles, so it is a nonzero regular form.
It is invariant because it was pulled back from \(Q_i\).  Equation
\eqref{eq:index-block-sections} now implies \(\chi_i^{m_*}=1\).

\smallskip\noindent\textbf{Combining the factor characters.}
The wedge of the rational pullbacks of the \(\omega_i\) to \(R\) is
the pullback of the product volume on \(P\).  It therefore generates
\(K_R\), even though \(R\) need not itself split as a product.  Its
character on \(G_0\) is the product of the block characters, and is
killed by \(m_*\).  A further factor depending only on
\([G:G_0]\le4!\) kills the character on all of \(G\).  For example,
one may multiply by the least common multiple of the integers from
one to \(4!\).  Equation \eqref{eq:index-character} proves the desired
uniform canonical multiple in this case.

\smallskip\noindent\textbf{The purely abelian case.}
If \(P=A\) is an abelian fourfold, take the same normal Galois closure
\(R\to X\).  The quasi-\'etale map \(R\to A\) is \'etale by purity.
A connected finite \'etale cover of a complex abelian variety is again
abelian: analytically it is obtained from a finite-index sublattice of
the lattice in its universal cover, and it is projective here.  Each
element of \(G\) acts integrally on
\(H^4(R,\Z)\), a free group of rank \(\binom84=70\).  Its eigenvalue on
the one-dimensional space \(H^{4,0}(R)\) is the canonical character.
If that eigenvalue has order \(d\), its cyclotomic polynomial divides an
integral characteristic polynomial of degree 70, and hence
\(\varphi(d)\le70\).  There are only finitely many such \(d\).  Their
least common multiple kills every value of the character, and
\eqref{eq:index-character} again gives a uniform trivial multiple of
\(K_X\).  Taking a common multiple for the two cases proves the
proposition.
\end{proof}

\subsection{Rational images of a single nonflat factor}

\begin{lemma}\label{lem:index-single-factor-images}
Suppose that a projective klt fourfold \(X\), with \(K_X\sim_{\Q}0\),
has a finite quasi-\'etale cover \(P\to X\) in which \(P\) is one
four-dimensional irreducible Calabi--Yau or irreducible holomorphic
symplectic factor as above.  Then the following statements hold.
\begin{enumerate}[label=\textup{(\roman*)}]
\item There is no dominant rational map from \(X\) to a positive
dimensional non-uniruled smooth projective variety of dimension less
than four.
\item Every positive dimensional rational image of dimension less
than four has rationally connected smooth projective resolution.
\item Every smooth projective model of \(X\) has irregularity zero.
\end{enumerate}
\end{lemma}

\begin{proof}
Suppose that \(X\dashrightarrow T\) contradicts (i), and set
\(j=\dim T\in\{1,2,3\}\).  Non-uniruledness implies that \(K_T\) is
pseudo-effective by \cite{BDPP2013}.  The minimal-model and abundance
theorems in dimension at most three imply \(\kappa(T,K_T)\ge0\).
In dimension three, one takes a projective \(\Q\)-factorial terminal
minimal model and applies abundance for minimal threefolds
\cite{FlipsAbundance1992,Kawamata1992}; canonical sections are unchanged
on a resolution.  The lower dimensional cases are the corresponding
curve and surface statements, or follow by first taking a product with
an abelian variety of dimension \(3-j\) and applying the same
three-dimensional results.

Choose a nonzero section \(\eta\in H^0(T,mK_T)\).  Resolve the composite
rational map from \(P\) to \(T\), obtaining a smooth projective
\(u:W\to P\) and a morphism \(f:W\to T\).  Pullback by differentials
gives a nonzero tensor
\begin{equation}\label{eq:index-pulled-tensor}
         s=f^*\eta\in H^0\bigl(W,(\Omega_W^j)^{\otimes m}\bigr).
\end{equation}
It is nonzero because the map is dominant in characteristic zero.
At the function-field level it has the form
\(s=a\beta^{\otimes m}\), where \(a\in\C(P)^*\) and \(\beta\) is a
nonzero decomposable rational \(j\)-form, obtained by pulling back a
rational top form on \(T\).

We claim that the divisorial zero locus of \(s\) is \(u\)-exceptional.
Choose an ample Cartier divisor \(H_P\) on \(P\), and put
\(H=u^*H_P\).  Since \(P\) is canonical with trivial canonical class,
\(K_W\) is \(\Q\)-linearly equivalent to an effective exceptional
divisor.  Thus it is pseudo-effective and \(K_W\cdot H^3=0\).
Generic semipositivity gives nonnegative degrees for torsion-free
quotients of \(\Omega_W^1\) with respect to any ample polarization
\cite[Theorem~2.1]{CampanaPaun2015}.  It gives the same conclusion for
torsion-free quotients of
\(\mathcal V=(\Omega_W^j)^{\otimes m}\).

Here is the tensor step explicitly.  Use an ample rational perturbation
\(H+\epsilon H_0\), where \(H_0\) is ample and \(\epsilon>0\).
The Harder--Narasimhan factors of \(\Omega_W^1\) have nonnegative slopes.
On a sufficiently high general complete-intersection curve the
filtration restricts to a filtration by semistable bundles
\cite[Theorem~6.1]{MehtaRamanathan1982}; the curve can avoid all the
codimension-two loci where the sheaves or quotients fail to be locally
free.  Tensor products of semistable bundles in characteristic zero
are semistable, with the sum of the slopes
\cite[Theorem~3.18]{RamananRamanathan1984}.  Tensor powers consequently
have filtrations whose slopes are nonnegative.  Exterior powers are
quotients of tensor powers, so the same lower slope bound applies to
\(\mathcal V\) and to its torsion-free quotients.  Letting
\(\epsilon\) decrease to zero proves the required nonnegative degree
against \(H^3\).

Saturate the rank-one subsheaf generated by \(s\) in \(\mathcal V\).
It is a line bundle \(\OO_W(Z)\), with \(Z\ge0\) the divisorial zero
divisor of \(s\), and the quotient is torsion-free.  The first Chern
class of \(\mathcal V\) is a positive integral multiple of \(K_W\).
It has degree zero against \(H^3\); the preceding quotient inequality
therefore gives
\[
                              Z\cdot H^3\le0.
\]
Every nonexceptional prime divisor has strictly positive intersection
with \(H^3\), by the projection formula and ampleness on \(P\).
Since \(Z\) is effective, it follows that \(Z\) is exceptional, as
claimed.

At the generic point of a prime divisor on \(P\), write the rational
form \(\beta\) in a regular frame of \(\Omega_P^j\), and let \(b\in\Z\)
be the minimum order of its coefficients.  The absence of divisorial
zeros of \(s\) at that point says
\[
                            \ord(a)+mb=0.
\]
Hence every coefficient of \(\Div_P(a)\) is divisible by \(m\).
Normalize \(P\) in a field component of the extension obtained by
adjoining \(z\) with \(z^m=a\).  This cover is finite and quasi-\'etale:
at any divisorial discrete valuation, removing the uniformizer to its
\(m\)-divisible order reduces the equation to a root of a unit, which
is unramified in characteristic zero.  The rational form \(z\beta\)
has no divisorial poles.  It is therefore a nonzero reflexive
\(j\)-form on the cover, and remains decomposable at its generic point.

This contradicts the defining form algebra of \(P\).  For a
four-dimensional irreducible Calabi--Yau factor there are no such
forms in degrees one, two or three on any quasi-\'etale cover.  For a
four-dimensional irreducible symplectic factor the only possible
degree is two, and every such form is a scalar multiple of the
generically nondegenerate symplectic form.  Its square is nonzero,
whereas the square of a decomposable two-form is zero.  This proves (i).

For (ii), resolve the rational image and take its maximal rationally
connected fibration.  Its smooth projective base is non-uniruled
\cite[Corollary~1.4]{GHS2003}.  A positive dimensional such base would
be a rational image of \(X\) of dimension less than four, contrary to
(i).  The base must be a point, which is exactly rational
connectedness of the resolution.

Finally let \(W_X\) be a smooth projective model of \(X\).  A regular
one-form on \(W_X\) pulls back under the dominant generically finite
rational map \(P\dashrightarrow W_X\) to a reflexive one-form on
\(P\).  Indeed, the rational map to the proper target is defined at
all codimension-one points, and the pulled-back form is regular there;
reflexive extension completes it.  Pullback is injective at the
function-field level in characteristic zero.  The absence of
one-forms on \(P\) forces the original form to vanish.  Hodge
decomposition and symmetry on the smooth projective \(W_X\) then give
\(h^1(W_X,\OO_{W_X})=h^0(W_X,\Omega^1_{W_X})=0\), proving (iii).
\end{proof}

We also need the following numerical fact about exceptional divisors.
Its constants depend on the fixed birational morphism only; no uniform
birational boundedness is involved.

\begin{lemma}\label{lem:index-exceptional-psef}
Let \(u:W\to X\) be a projective birational morphism from a smooth
projective fourfold to a normal projective variety.  A pseudo-effective
\(u\)-exceptional real divisor on \(W\) is effective.
\end{lemma}

\begin{proof}
Fix a very ample Cartier divisor \(L\) on \(X\), put \(A=u^*L\), and
fix an ample Cartier divisor \(H\) on \(W\).  Because \(A\) is big,
there are an integer \(b>0\) and an effective Cartier divisor \(F\)
such that \(bA\sim H+F\).  If \(S\) is a nonexceptional prime divisor
not contained in \(\Supp F\), then \(F\cdot S\) is an effective
codimension-two cycle.  Both \(bA\) and \(H\) are nef, and factoring
their cubes gives
\[
 \begin{split}
 (bA)^3\cdot S-H^3\cdot S
   &=F\cdot S\cdot\bigl((bA)^2+bA\,H+H^2\bigr)\ge0.
 \end{split}
\]
There are only finitely many nonexceptional primes in \(\Supp F\), and
each has positive \(A^3\)-degree.  Increasing a constant to handle these
primes gives
\begin{equation}\label{eq:index-strict-degree}
              H^3\cdot S\le C A^3\cdot S
                    =C L^3\cdot u_*S
\end{equation}
for every nonexceptional prime \(S\).

Let \(E=\sum_i a_iE_i\) be exceptional and pseudo-effective.  Choose
effective real divisors \(D_n\) whose numerical classes converge to
\([E]\).  Write \(D_n=S_n+T_n\), where \(S_n\) has no exceptional
components and \(T_n\) is supported on the finitely many exceptional
prime divisors of \(u\).  Projection gives
\[
      A^3\cdot S_n=A^3\cdot D_n\longrightarrow A^3\cdot E=0.
\]
Equation \eqref{eq:index-strict-degree} implies \(H^3\cdot S_n\to0\).
Choose divisor classes \(Q_1,\ldots,Q_\rho\) forming a basis of
\(N^1(W)_{\R}\), and choose \(c_j>0\) so that both
\(c_jH+Q_j\) and \(c_jH-Q_j\) are ample.  Effectivity gives
\[
       |Q_j\cdot H^2\cdot S_n|\le c_j H^3\cdot S_n\longrightarrow0.
\]
The Hodge index theorem makes the form
\((D,Q)\mapsto D\cdot Q\cdot H^2\) nondegenerate on
\(N^1(W)_{\R}\).  Thus \([S_n]\to0\), and consequently
\([T_n]\to[E]\).

The numerical classes of the exceptional prime divisors are linearly
independent.  In fact, an exceptional numerically trivial linear
combination and its negative are both relatively nef.  The negativity
lemma forces both to be nonpositive, hence forces the combination to
vanish as a divisor \cite[Lemma~3.39]{KollarMori1998}.  Their positive
cone is therefore closed and has unique coordinates in these classes.
Since all \(T_n\) belong to that cone, so does \([E]\); independence
then gives \(a_i\ge0\) for every \(i\).
\end{proof}

\begin{corollary}\label{cor:index-noncanonical}
The canonical orders of the noncanonical varieties in
Lemma~\ref{lem:index-single-factor-images} are uniformly bounded.
\end{corollary}

\begin{proof}
Choose a smooth projective resolution \(u:W\to X\) exhibiting a
negative ordinary discrepancy.  Its canonical divisor satisfies
\[
                  K_W=u^*K_X+E\sim_{\Q}E,
\]
where \(E\) is exceptional with a negative coefficient.  By
Lemma~\ref{lem:index-exceptional-psef}, \(K_W\) is not pseudo-effective.
The uniruledness criterion of \cite{BDPP2013} implies that \(W\) is
uniruled.  Its maximal rationally connected quotient consequently has
dimension less than four, and its base is non-uniruled
\cite[Corollary~1.4]{GHS2003}.  If that base had positive dimension, the
induced rational map from \(X\) would contradict
Lemma~\ref{lem:index-single-factor-images}(i).  Thus \(W\) is rationally
connected.  Apply Proposition~\ref{prop:rc-torsion} directly to
\((X,0)\), with zero nef data.  Its hypotheses are precisely klt
singularities, \(K_X\sim_{\Q}0\), and a rationally connected smooth
resolution, all of which have been checked.  Its uniform integer
bounds the canonical order.
\end{proof}

\subsection{Terminal models and intermediate fibrations}

\begin{lemma}\label{lem:index-terminalization}
A normal projective canonical variety \(X\) with \(K_X\sim_{\Q}0\)
has a projective crepant \(\Q\)-factorial terminalization \(V\to X\).
The canonical orders of \(V\) and \(X\) are equal.
\end{lemma}

\begin{proof}
On a fixed smooth resolution, the discrepancy divisor of \(X\) is
effective.  Any divisor exceptional over that resolution has positive
ordinary discrepancy over the smooth resolution, and adding the
pullback of an effective discrepancy divisor preserves positivity.
Consequently only the finitely many exceptional divisors already on
the fixed resolution can have ordinary discrepancy zero over \(X\).

The projective extraction theorem
\cite[Corollary~1.4.3]{BCHM2010} extracts this finite list and gives a
\(\Q\)-factorial variety.  One may arrange its klt application as
follows.  Fix an integer \(r\) with \(rK_X\) Cartier, and choose an
effective ample Cartier divisor \(H\) containing all the centers of
the selected valuations.  Since \(X\) is klt, choose \(c>0\) such that
\((X,cH)\) is lc.  Then, for every divisorial valuation \(E\),
\[
                      \ord_E(H)\le A_X(E)/c.
\]
Every exceptional log discrepancy greater than one is at least
\(1+1/r\).  Choosing a positive rational \(\delta\) with
\(\delta/c<1/(r+1)\) keeps all those discrepancies greater than one,
while the selected log discrepancies equal to one decrease strictly
below one.  Take \(\delta\) still smaller if necessary so that
\((X,\delta H)\) is klt, and apply the extraction theorem to this
pair and exactly the selected list.  The list can be empty, in which
case the theorem gives a small \(\Q\)-factorialization.

Restore boundary zero.  The extracted divisors have ordinary
discrepancy zero for \(X\), so the resulting morphism is crepant.
All divisors still exceptional over the extracted variety have
positive ordinary discrepancy; hence the variety is terminal.
Finally, compatible canonical divisors satisfy \(K_V=u^*K_X\).
A principal multiple downstairs pulls back to the same principal
multiple upstairs, and a principal multiple upstairs pushes forward
to the same one downstairs.  This proves equality of their exact
orders.
\end{proof}

\begin{lemma}\label{lem:index-moving-divisor}
There is a uniform positive integer \(N_{\rm mov}\) with the following
property.  Let \(V\) be a projective \(\Q\)-factorial terminal fourfold
with \(K_V\sim_{\Q}0\), and suppose that every positive dimensional
rational image of dimension less than four has rationally connected
smooth projective resolution.  If \(V\) has an effective Weil divisor
\(D\) with \(1\le\kappa(V,D)\le3\), then
\(N_{\rm mov}K_V\sim0\).
\end{lemma}

\begin{proof}
By \(\Q\)-factoriality, \(D\) is \(\Q\)-Cartier.  Choose a small
positive rational number \(\delta\) for which \((V,\delta D)\) is klt.
Its adjoint has Iitaka dimension \(\kappa(V,D)\), since
\(K_V\sim_{\Q}0\).  Apply Proposition~\ref{prop:semiample-model},
starting its birational program at \(V\), to obtain a \(\Q\)-factorial
klt model \(V'\) and its semiample contraction
\[
                       f:V'\longrightarrow Z,
          \qquad 1\le\dim Z=\kappa(V,D)\le3.
\]

We check the canonical divisor separately from the perturbed adjoint.
Choose compatible canonical representatives and a rational function
\(h\) with \(rK_V=\Div(h)\).  Each birational contraction in the program
pushes this same relation to the target, and each flip carries it to
the new model by its codimension-one identification.  All these
models are \(\Q\)-factorial.  On a common resolution their canonical
pullbacks are therefore the same divisor \(r^{-1}\Div(h)\).
In particular, \(K_{V'}\sim_{\Q}0\), the zero-boundary varieties are
crepant, and their canonical orders are equal.  They remain klt;
alternatively, this follows from crepant equality starting at the
terminal \(V\).

Now use Theorem~\ref{thm:relative-denominators} on the contraction
\(f:(V',0)\to Z\), with zero boundary and base class zero.  It gives
uniform integers \(p_0\mid p\), a rational function \(\psi\), and
generalized klt data satisfying
\begin{equation}\label{eq:index-zero-relative}
 K_{V'}+\frac1{p_0}\Div(\psi)
       =f^*D_Z,\qquad
 D_Z=K_Z+B_Z+M_Z\sim_{\Q}0,
\end{equation}
where \(B_Z\ge0\) has coefficients in a fixed DCC set and the nef
b-divisor has Cartier denominator dividing \(p\).  The base \(Z\) is a
rational image of \(V\), so its smooth projective resolution is
rationally connected.  Proposition~\ref{prop:rc-torsion} gives a
uniform \(l\), which we enlarge to a multiple of \(p_0\), such that
\(lD_Z\sim0\).  Multiplying \eqref{eq:index-zero-relative} by \(l\)
then gives \(lK_{V'}\sim0\), hence \(lK_V\sim0\).

The coefficient \(\delta\) and the divisor \(D\) were used only to
produce a semiample model.  The two uniform statements were applied
afterward with boundary exactly zero, in one of the three fixed base
dimensions.  Their integers thus depend on neither \(D\) nor
\(\delta\).  A common multiple for these dimensions is
\(N_{\rm mov}\).
\end{proof}

\begin{lemma}\label{lem:index-general-type-fibres}
Let \(V\) be a projective \(\Q\)-factorial terminal fourfold with
\(K_V\sim_{\Q}0\).  Suppose that every effective Weil divisor on
\(V\) of positive Iitaka dimension is big.  Let
\(\pi:Y\to V\) be its cyclic index cover, with deck group \(G\).
Then every \(G\)-equivariant rational fibration of \(Y\) with positive
dimensional base and fibre has a geometric generic fibre of general
type on a smooth resolution.  The same conclusion holds for every
such rational fibration of \(V\).
\end{lemma}

\begin{proof}
The variety \(Y\) is terminal.  In fact, an exceptional divisorial
valuation over \(Y\) restricts to an exceptional valuation over \(V\),
and the finite discrepancy formula multiplies its log discrepancy
by its positive ramification index.  The log discrepancy remains
strictly greater than one.  Also \(K_Y\sim0\).

Consider a \(G\)-equivariant rational fibration \(Y\dashrightarrow T\).
Its base field \(\C(T)\) is relatively algebraically closed in
\(\C(Y)\).  The invariant field \(\C(T)^G\) lies in
\(\C(Y)^G=\C(V)\), has the same positive transcendence degree as
\(\C(T)\), and is relatively algebraically closed in \(\C(V)\).
For the last assertion, an element of \(\C(V)\) algebraic over
\(\C(T)^G\) is algebraic over \(\C(T)\), hence belongs to
\(\C(T)\), and then is invariant.  A projective model of the invariant
field therefore gives a descended dominant rational map
\(V\dashrightarrow Z\), where \(0<\dim Z<4\).

Fix smooth projective resolutions of the maps and the bases.  Thus
we may use a smooth projective birational morphism \(u:U\to Y\),
a morphism \(f:U\to T'\) to a smooth model of the unquotiented base,
and a morphism \(g:U\to Z'\) to a projective model of the quotient
base.  The map \(g\) is constant on the generic fibre of \(f\).
Take a general hyperplane divisor \(H\) for a very ample embedding
of \(Z'\).  On the domain of \(V\dashrightarrow Z'\), pull it back
and take its closure to obtain an effective Weil divisor \(D\) on
\(V\).  The domain contains every codimension-one point, because
\(V\) is normal and \(Z'\) is proper.  The hyperplane pulls form a
moving system of linearly equivalent divisors: the rational-function
ratios on the domain extend as the same principal-divisor relations
on \(V\).  Some ratio is nonconstant, since \(\dim Z'>0\).  Thus
\(\kappa(V,D)\ge1\), and \(D\) is big by hypothesis.

Write \(q=\pi\circ u\).  We have
\begin{equation}\label{eq:index-exceptional-fibre}
                           q^*D=g^*H+E,
\end{equation}
where \(E\) is an effective \(u\)-exceptional rational divisor.
The pullback exists because \(V\) is \(\Q\)-factorial.  At every
nonexceptional generic divisorial point of \(U\), the two terms
\(q^*D\) and \(g^*H\) agree: the finite cover takes that point to a
codimension-one point of \(V\), where the original rational map is
defined.  Thus their difference is exceptional.  After the resolution
has been fixed, choose \(H\) general enough that it contains none of
the generic images under \(g\) of its finitely many exceptional
prime divisors.  The coefficient of \(g^*H\) at each such prime is
then zero, while \(q^*D\) is effective.  This proves the claimed
effectivity of \(E\).

Let \(F\) be the smooth geometric generic fibre of \(f\).  The divisor
\(q^*D\) is big on \(U\), and its restriction to \(F\) is big as
well.  Indeed, write a big rational class as the sum of an ample
rational class and an effective rational class.  The integral
geometric generic fibre is not contained in the support of the
effective term, and the ample term restricts amply.  The class
\(g^*H\) restricts trivially to \(F\), so
\eqref{eq:index-exceptional-fibre} implies that \(E|_F\) is big.

Terminality and \(K_Y\sim0\) give
\[
                        K_U\sim_{\Q}\sum_i a_iE_i,
                       \qquad a_i>0,
\]
where the sum runs over all \(u\)-exceptional prime divisors.
Since \(E\) is effective and supported on this finite set, there is
a positive rational number \(c\) such that
\(\sum_i a_iE_i-cE\ge0\).  Restriction to the geometric generic
fibre is effective.  Smooth-fibre adjunction identifies
\(K_U|_F\) with \(K_F\), and therefore \(K_F\) is the sum, up to
\(\Q\)-linear equivalence, of the big divisor \(cE|_F\) and an
effective divisor.  It is big, which proves that \(F\) is of
general type.

For a rational fibration of \(V\), use the identical argument with
\(Y=V\), \(\pi=\id\), and the trivial group.  Terminality of \(V\)
provides the same positive discrepancy comparison.
\end{proof}


We have now removed the product, noncanonical and intermediate-divisor
cases by uniform index bounds. In a hypothetical sequence of unbounded
indices, its remaining terminal models $V$ therefore have irregularity
zero and no effective nonbig divisor of positive Iitaka dimension.
Lemma~\ref{lem:index-general-type-fibres} gives the requisite general-type
fibres for both $V$ and its canonical cover $Y$. The next three sections
establish chain degree bounds and the scalar and product curve estimates
for these varieties;
Section~\ref{sec:transports} will turn their discrepancy into birational
self-maps, and Section~\ref{sec:index-zero} will finish the contradiction.
